\chapter{从作用量到传播与辐射}

\begin{prereqbox}
熟悉四维指标记号和分部积分。采用度规
$g_{\mu\nu}=\mathrm{diag}(1,-1,-1,-1)$，
$x^\mu=(ct,\vb r)$，
$A^\mu=(\phi/c,\vb A)$，
$J^\mu=(c\rho,\vb J)$。重复指标求和。不同教材的度规和势定义会改变若干符号，面试时应先声明约定。
\end{prereqbox}

\section{为什么从作用量出发}
Maxwell 方程若只作为四条经验公式记忆，规范自由度、守恒律和相对论结构彼此分离。作用量把它们统一：
\[
S[A]=\int\dd^4x\,\mathcal L,\qquad
\mathcal L=-\frac1{4\mu_0}F_{\mu\nu}F^{\mu\nu}-J_\mu A^\mu,
\]
\[
F_{\mu\nu}=\partial_\mu A_\nu-\partial_\nu A_\mu.
\]
第一项是自由场，第二项是外源与场的最小线性耦合。系数 $1/4$ 同时补偿
$F_{\mu\nu}$ 反对称导致的重复计数和变分时出现的因子 $2$。

\section{逐项变分推出非齐次 Maxwell 方程}
令 $A_\mu\to A_\mu+\delta A_\mu$，并要求变分在时空积分边界消失。首先
\[
\delta F_{\mu\nu}
=\partial_\mu\delta A_\nu-\partial_\nu\delta A_\mu.
\]
自由场项的变分是
\[
\delta S_f=-\frac1{4\mu_0}\int\dd^4x\,
\delta(F_{\mu\nu}F^{\mu\nu})
=-\frac1{2\mu_0}\int\dd^4x\,F^{\mu\nu}\delta F_{\mu\nu}.
\]
代入 $\delta F_{\mu\nu}$：
\[
\delta S_f=-\frac1{2\mu_0}\int\dd^4x\,
F^{\mu\nu}(\partial_\mu\delta A_\nu-\partial_\nu\delta A_\mu).
\]
第二项交换哑指标 $\mu\leftrightarrow\nu$，并用
$F^{\nu\mu}=-F^{\mu\nu}$：
\[
-F^{\mu\nu}\partial_\nu\delta A_\mu
=+F^{\nu\mu}\partial_\mu\delta A_\nu
=-F^{\mu\nu}\partial_\mu\delta A_\nu.
\]
因此两项相同，
\[
\delta S_f=-\frac1{\mu_0}\int\dd^4x\,
F^{\mu\nu}\partial_\mu\delta A_\nu.
\]
对 $\partial_\mu$ 分部积分：
\[
\delta S_f=
-\frac1{\mu_0}\int_{\partial\Omega}\dd\Sigma_\mu
F^{\mu\nu}\delta A_\nu
+\frac1{\mu_0}\int_\Omega\dd^4x\,
(\partial_\mu F^{\mu\nu})\delta A_\nu.
\]
边界变分为零，第一项消失。源项给
\[
\delta S_J=-\int\dd^4x\,J^\nu\delta A_\nu.
\]
合并：
\[
\delta S=\int\dd^4x\left[
\frac1{\mu_0}\partial_\mu F^{\mu\nu}-J^\nu
\right]\delta A_\nu.
\]
由于内部任意 $\delta A_\nu$ 都应使 $\delta S=0$，得到
\[
\boxed{\partial_\mu F^{\mu\nu}=\mu_0J^\nu}.
\]

\subsection{拆回三维形式}
取 $\nu=0$，利用
$F^{i0}=E_i/c$ 与 $\partial_i=\partial/\partial x_i$：
\[
\partial_iF^{i0}=\frac1c\nabla\cdot\vb E
=\mu_0c\rho,
\]
即
\[
\nabla\cdot\vb E=\mu_0c^2\rho=\rho/\varepsilon_0.
\]
取 $\nu=i$，时间项和空间项分别产生位移电流与磁场旋度，得到
\[
\nabla\times\vb B-\frac1{c^2}\pdv{\vb E}{t}
=\mu_0\vb J.
\]
这正是 Gauss 电定律和 Ampère--Maxwell 定律。

\section{齐次方程不是新的运动方程}
因为 $F=\dd A$，恒有
\[
\partial_\lambda F_{\mu\nu}
+\partial_\mu F_{\nu\lambda}
+\partial_\nu F_{\lambda\mu}=0.
\]
把 $F_{\mu\nu}=\partial_\mu A_\nu-\partial_\nu A_\mu$ 展开，所有二阶混合偏导两两抵消。这是 Bianchi 恒等式，不依赖作用量极值。其三维内容为
\[
\boxed{\nabla\cdot\vb B=0,\qquad
\nabla\times\vb E=-\pdv{\vb B}{t}}.
\]
因此四条 Maxwell 方程中，两条是作用量给出的动力学方程，两条来自用势表示场的几何恒等式。

\section{规范不变性与电荷守恒}
规范变换
\[
A_\mu\to A_\mu+\partial_\mu\chi
\]
不改变 $F_{\mu\nu}$。源耦合的变化为
\[
\delta_\chi S_J=-\int\dd^4x\,J^\mu\partial_\mu\chi.
\]
分部积分：
\[
\delta_\chi S_J
=-\int_{\partial\Omega}\dd\Sigma_\mu J^\mu\chi
+\int_\Omega\dd^4x\,\chi\,\partial_\mu J^\mu.
\]
边界项消失时，作用量对任意 $\chi$ 规范不变要求
\[
\boxed{\partial_\mu J^\mu=0}
\quad\Longleftrightarrow\quad
\pdv\rho t+\nabla\cdot\vb J=0.
\]
反过来，对场方程取 $\partial_\nu$：
\[
\partial_\nu\partial_\mu F^{\mu\nu}
=\mu_0\partial_\nu J^\nu.
\]
左边因偏导对称、$F$ 反对称而为零，也推出连续性方程。

\begin{pitfallbox}
规范变换是同一物理场的冗余描述，不是把一种物理状态变成另一种。所谓“规范对称性自发破缺”在严格表述中需要谨慎；超导中真正可观测的是规范不变相位梯度、磁通量子和相位刚度。
\end{pitfallbox}

\section{Lorenz 规范与波动方程}
将 $F^{\mu\nu}$ 展开：
\[
\partial_\mu F^{\mu\nu}
=\partial_\mu\partial^\mu A^\nu
-\partial^\nu(\partial_\mu A^\mu).
\]
定义 d'Alembert 算符
\[
\Box=\partial_\mu\partial^\mu
=\frac1{c^2}\frac{\partial^2}{\partial t^2}-\nabla^2.
\]
选择 Lorenz 规范
\[
\partial_\mu A^\mu=0
\quad\Longleftrightarrow\quad
\nabla\cdot\vb A+\frac1{c^2}\pdv\phi t=0,
\]
场方程化为四个波动方程
\[
\boxed{\Box A^\nu=\mu_0J^\nu}.
\]
规范并未完全固定：若 $\Box\chi=0$，则
$A_\mu+\partial_\mu\chi$ 仍满足 Lorenz 规范，这称剩余规范自由度。

\section{波动方程的推迟 Green 函数}
要求 Green 函数满足
\[
\Box G(x-x')=\delta^{(4)}(x-x').
\]
解可写为
\[
A^\mu(x)=\mu_0\int\dd^4x'\,G(x-x')J^\mu(x').
\]
物理上要求原因先于结果，即 $G(t<t')=0$，选择推迟 Green 函数。三维形式为
\[
G_{\rm ret}(\vb r,t;\vb r',t')
=\frac{\delta(t-t'-R/c)}{4\pi R},
\qquad R=|\vb r-\vb r'|.
\]
把它代回并先做 $t'$ 积分：
\[
\boxed{
\phi(\vb r,t)=\frac1{4\pi\varepsilon_0}
\int\frac{\rho(\vb r',t-R/c)}R\,\dd^3r'},
\]
\[
\boxed{
\vb A(\vb r,t)=\frac{\mu_0}{4\pi}
\int\frac{\vb J(\vb r',t-R/c)}R\,\dd^3r'}.
\]
源必须在推迟时刻 $t_r=t-R/c$ 取值。传播速度 $c$ 并非事后加入，而是双曲型场方程和因果边界条件的结果。

\subsection{验证 Green 函数结构}
Fourier 变换到 $(\vb k,\omega)$：
\[
G(\vb k,\omega)
=\frac1{k^2-(\omega/c)^2}.
\]
极点怎样绕过实轴决定边界条件。推迟解对应把极点移到保证
$t<t'$ 时积分闭合后为零的一侧。静态极限 $\omega=0$ 给
$G(\vb k)=1/k^2$，反变换就是 $1/4\pi R$ 的 Coulomb Green 函数。

\section{静电多极展开}
局域电荷分布外部 $r\gg r'$：
\[
\frac1{|\vb r-\vb r'|}
=\frac1r+\frac{\vb r\cdot\vb r'}{r^3}
+\frac{3(\vb r\cdot\vb r')^2-r^2r'^2}{2r^5}
+\cdots.
\]
代入
$\phi=(4\pi\varepsilon_0)^{-1}\int\rho(\vb r')/|\vb r-\vb r'|\dd^3r'$。
定义
\[
Q=\int\rho\,\dd^3r',\qquad
\vb p=\int\vb r'\rho\,\dd^3r',
\]
\[
Q_{ij}=\int\rho(3x_i'x_j'-r'^2\delta_{ij})\dd^3r'.
\]
于是
\[
\phi(\vb r)=\frac1{4\pi\varepsilon_0}
\left[\frac Qr+\frac{\vb p\cdot\hat{\vb r}}{r^2}
+\frac1{2r^3}Q_{ij}\hat r_i\hat r_j+\cdots\right].
\]
若总电荷 $Q\ne0$，偶极矩依赖原点；若 $Q=0$，平移原点不改变 $\vb p$。远场最低非零多极矩支配势。

\section{唯一性定理与镜像法}
若两组势 $\phi_1,\phi_2$ 在区域内满足同一 Poisson 方程并具有相同 Dirichlet 边界值，差
$u=\phi_1-\phi_2$ 满足 $\nabla^2u=0$ 且边界 $u=0$。利用 Green 第一恒等式：
\[
\int_V|\nabla u|^2\dd V
=\oint_{\partial V}u\pdv u n\dd a-\int_Vu\nabla^2u\dd V=0.
\]
所以 $\nabla u=0$，$u$ 为常数；边界令常数为零。镜像法只需在物理区域内满足 Poisson 方程和导体边界条件，唯一性定理便保证解正确。

\section{Poynting 定理逐步推导}
从 Ampère--Maxwell 方程点乘 $\vb E$：
\[
\vb E\cdot(\nabla\times\vb H)
=\vb J\cdot\vb E+\vb E\cdot\pdv{\vb D}{t}.
\]
从 Faraday 方程点乘 $\vb H$：
\[
\vb H\cdot(\nabla\times\vb E)
=-\vb H\cdot\pdv{\vb B}{t}.
\]
用恒等式
\[
\nabla\cdot(\vb E\times\vb H)
=\vb H\cdot\nabla\times\vb E
-\vb E\cdot\nabla\times\vb H
\]
相减，在线性无色散介质中
\[
\vb E\cdot\pdv{\vb D}{t}
=\pdv{}{t}\left(\frac12\vb E\cdot\vb D\right),
\]
\[
\vb H\cdot\pdv{\vb B}{t}
=\pdv{}{t}\left(\frac12\vb B\cdot\vb H\right).
\]
得到
\[
\boxed{\pdv u t+\nabla\cdot\vb S=-\vb J\cdot\vb E},
\]
\[
u=\frac12(\vb E\cdot\vb D+\vb B\cdot\vb H),
\qquad \vb S=\vb E\times\vb H.
\]
体积分形式表明：区域内场能变化 $+$ 流出表面的能流 $=$ 场对物质做功的负值。

\section{平面波、介质与 Fresnel 系数}
无源均匀介质取
$\vb E=\vb E_0e^{i(\vb k\cdot\vb r-\omega t)}$。
Maxwell 方程给
\[
\vb k\cdot\vb E_0=0,\qquad
\vb B_0=\frac1\omega\vb k\times\vb E_0,
\]
\[
k^2=\mu\varepsilon\omega^2,\qquad
v=\frac1{\sqrt{\mu\varepsilon}}.
\]
法向入射界面，切向 $E,H$ 连续。设入射、反射、透射振幅
$E_i,E_r,E_t$，波阻抗 $Z=\sqrt{\mu/\varepsilon}$：
\[
E_i+E_r=E_t,\qquad
\frac{E_i-E_r}{Z_1}=\frac{E_t}{Z_2}.
\]
解得
\[
r=\frac{E_r}{E_i}=\frac{Z_2-Z_1}{Z_2+Z_1},
\qquad
t=\frac{E_t}{E_i}=\frac{2Z_2}{Z_1+Z_2}.
\]
功率反射率 $R=|r|^2$；无耗散时需用 Poynting 流验证
$R+T=1$，其中 $T$ 含阻抗比。

\section{矩形波导的截止频率}
理想导体内切向电场为零。设波导横截面 $0<x<a,0<y<b$，沿
$z$ 传播。对 TM 模，纵向电场满足
\[
(\partial_x^2+\partial_y^2+k_c^2)E_z=0,
\quad E_z|_{\partial A}=0.
\]
分离变量得到
\[
E_z=E_0\sin\frac{m\pi x}{a}
\sin\frac{n\pi y}{b},
\]
\[
k_c^2=\left(\frac{m\pi}{a}\right)^2+
\left(\frac{n\pi}{b}\right)^2.
\]
纵向传播常数
\[
\beta^2=\frac{\omega^2}{c^2}-k_c^2.
\]
当 $\omega<ck_c$，$\beta$ 为虚数，模式指数衰减；截止频率
\[
\omega_c=ck_c.
\]
波导中场一般含纵向分量，因此“电磁波必为横波”只对均匀无界介质平面波成立。

\section{局域源的远区展开}
对尺寸 $d$ 的源，观察点 $r\gg d$：
\[
R=|\vb r-\vb r'|\simeq r-\hat{\vb r}\cdot\vb r',
\qquad
\frac1R\simeq\frac1r.
\]
推迟时间
\[
t-\frac Rc\simeq t-\frac rc+
\frac{\hat{\vb r}\cdot\vb r'}c.
\]
若源还满足长波条件 $d\ll\lambda$，最后一项相位可作展开。最低阶有时间变化的电偶极矩
\[
\vb p(t)=\int\vb r'\rho(\vb r',t)\dd^3r'.
\]
利用连续性方程：
\[
\dot{\vb p}
=\int\vb r'\pdv\rho t\dd^3r'
=-\int\vb r'\nabla'\cdot\vb J\,\dd^3r'.
\]
分部积分且局域源边界项消失：
\[
\boxed{\dot{\vb p}=\int\vb J\,\dd^3r'}.
\]
因此远区矢势最低阶为
\[
\vb A(\vb r,t)\simeq
\frac{\mu_0}{4\pi r}\dot{\vb p}(t-r/c).
\]

\section{电偶极辐射}
远区只保留 $1/r$ 项。由
$\vb B=\nabla\times\vb A$，对推迟时间求空间导数产生
$-\,\hat{\vb r}\partial_t/c$：
\[
\boxed{\vb B_{\rm rad}
=\frac{\mu_0}{4\pi cr}
\ddot{\vb p}(t-r/c)\times\hat{\vb r}}.
\]
远区局部像平面波，
\[
\boxed{\vb E_{\rm rad}
=c\,\vb B_{\rm rad}\times\hat{\vb r}
=\frac1{4\pi\varepsilon_0c^2r}
\hat{\vb r}\times(\ddot{\vb p}\times\hat{\vb r})}.
\]
若 $\vb p=p_0\cos\omega t\,\hat{\vb z}$，
\[
E_\theta=\frac{p_0\omega^2}{4\pi\varepsilon_0c^2r}
\sin\theta\cos[\omega(t-r/c)].
\]
周期平均 Poynting 流
\[
\avg S_r=\frac{p_0^2\omega^4}{32\pi^2
\varepsilon_0c^3r^2}\sin^2\theta.
\]
每立体角功率
\[
\boxed{\dv{\avg P}{\Omega}
=\frac{p_0^2\omega^4}{32\pi^2\varepsilon_0c^3}
\sin^2\theta}.
\]
积分
$\int\sin^2\theta\,\dd\Omega=8\pi/3$：
\[
\boxed{\avg P
=\frac{p_0^2\omega^4}{12\pi\varepsilon_0c^3}}.
\]
辐射沿偶极轴为零、赤道最大，且有强烈的 $\omega^4$ 依赖。

\section{近区、感应区和辐射区}
振荡偶极完整场包含
\[
E\sim\frac p{r^3}
+\frac{\dot p}{cr^2}
+\frac{\ddot p}{c^2r}.
\]
$1/r^3$ 是准静电近场，$1/r^2$ 是感应场，$1/r$ 是辐射场。只有辐射项使
$r^2S_r$ 在 $r\to\infty$ 保持有限总功率。近场能量在源和场之间往返交换，不代表净辐射。

\section{点电荷、Liénard--Wiechert 势与 Larmor 公式}
点电荷轨迹 $\vb r_q(t)$ 的源为
\[
\rho(\vb r,t)=q\delta[\vb r-\vb r_q(t)],
\qquad
\vb J=\rho\,\vb v(t).
\]
代入推迟势时，$\delta$ 函数的变量也依赖推迟时刻，积分产生 Jacobian
$\kappa=1-\hat{\vb n}\cdot\vb\beta$。得到
\[
\phi(\vb r,t)=
\left.\frac{q}{4\pi\varepsilon_0}
\frac1{\kappa R}\right|_{\rm ret},
\qquad
\vb A=\left.\frac{\vb v}{c^2}\phi\right|_{\rm ret}.
\]
求导后电场分成速度场与加速度场：
\[
\vb E=
\frac{q}{4\pi\varepsilon_0}
\left[
\frac{(1-\beta^2)(\hat{\vb n}-\vb\beta)}
{\kappa^3R^2}
+\frac{\hat{\vb n}\times[
(\hat{\vb n}-\vb\beta)\times\dot{\vb\beta}]}
c{\kappa^3R}
\right]_{\rm ret}.
\]
第一项 $\sim1/R^2$，第二项 $\sim1/R$ 并负责辐射。非相对论极限
$\beta\ll1$：
\[
\vb E_{\rm rad}
=\frac{q}{4\pi\varepsilon_0c^2R}
\hat{\vb n}\times(\hat{\vb n}\times\vb a).
\]
角积分得到 Larmor 功率
\[
\boxed{P=\frac{q^2a^2}{6\pi\varepsilon_0c^3}}.
\]
它与偶极辐射一致：单电荷偶极矩 $\vb p=q\vb r_q$，故
$\ddot{\vb p}=q\vb a$。

\section{完整例题：短偶极天线}
长度 $\ell\ll\lambda$ 的细导线电流近似
$I(t)=I_0\cos\omega t$。端点电荷满足
$\dot q=I$，偶极矩幅值
\[
p_0=q_0\ell=\frac{I_0\ell}{\omega}.
\]
代入偶极总辐射功率：
\[
\avg P=\frac{I_0^2\ell^2\omega^2}
{12\pi\varepsilon_0c^3}.
\]
若用峰值电流定义辐射电阻
$\avg P=I_0^2R_{\rm rad}/2$，则
\[
R_{\rm rad}=\frac{\ell^2\omega^2}
{6\pi\varepsilon_0c^3}
=80\pi^2\left(\frac{\ell}{\lambda}\right)^2\Omega.
\]
这展示了场论推导如何转化为可测的电路参数。

\section{面试表达与自测}
\begin{interviewbox}
\textbf{30 秒：}自由电磁场作用量由唯一的最低阶 Lorentz 标量
$F_{\mu\nu}F^{\mu\nu}$ 构成。对势变分给非齐次 Maxwell 方程；
$F=\dd A$ 自动给齐次方程；规范不变性要求电流守恒。Lorenz 规范把势化成波动方程，推迟 Green 函数保证因果传播，远区展开再给偶极辐射。

\textbf{黑板版：}至少能完整写出
$\delta(F^2)$、利用反对称性合并两项、分部积分、消去边界项和得到
$\partial_\mu F^{\mu\nu}=\mu_0J^\nu$ 的五步。
\end{interviewbox}
\begin{selftestbox}
\begin{enumerate}
\item 为什么场作用量系数是 $-1/4\mu_0$？
\item 变分时反对称性怎样把两项合并？
\item 齐次 Maxwell 方程为何是恒等式？
\item 规范不变性和电荷守恒的逻辑关系是什么？
\item 推迟与超前 Green 函数满足同一微分方程，物理上如何选择？
\item 多极展开的适用条件是什么？
\item 唯一性定理怎样保证镜像法？
\item 从哪一个矢量恒等式推出 Poynting 定理？
\item Fresnel 功率透射率为什么不总等于 $|t|^2$？
\item 波导截止的数学来源是什么？
\item 偶极远场中为什么出现 $\ddot p$？
\item 如何按 $1/r$ 幂次区分近场和辐射场？
\item Liénard--Wiechert 势中的 $\kappa$ 从哪里来？
\item Larmor 公式与偶极辐射如何相互验证？
\end{enumerate}
\end{selftestbox}
